\documentclass{article}
\usepackage{tikz-feynman}
\usepackage{titling}

%Bibliography
\usepackage[backend=biber,style=phys,sorting=nyt]{biblatex}
\addbibresource{refs.bib}
\appto{\bibfont}{\small}

\setlength{\droptitle}{-4cm} 

\title{PAP331 - Computational Methods in High Energy Physics}
\author{Shounak Naik}
\date{Exercise 1}

%borrowed feynmann layout from previous project
\begin{document}
\maketitle
\vspace{3cm}
\textbf{Solution to question 2:}
\vspace{0.7cm}
\begin{figure}[h]
    \centering
    \begin{minipage}{0.49\textwidth}
        \centering
        \begin{tikzpicture}[scale=1.5]
            \begin{feynman}
                \vertex at (-2.7,1.5) (A);
                \vertex at (-2.7,-1.5) (B);
                \vertex at (-0.2,1) (C);
                \vertex at (-1.2,-1) (D);
                \vertex at (1,1.5) (E);
                \vertex at (1,-1.5) (F);
                \diagram*{
                    (A) -- [fermion, edge label={\(e^+\)}] (C);
                    (B) -- [fermion, edge label'={\(e^-\)}] (D);
                    (C) -- [boson, edge label={\(\gamma\)}] (E);
                    (D) -- [boson, edge label={\(\gamma\)}] (F);
                    (D) -- [fermion, edge label'={\( \)}] (C);
                };
            \end{feynman}
        \end{tikzpicture} 
    \end{minipage}
    \begin{minipage}{0.49\textwidth}
        \centering
        \begin{tikzpicture}[scale=1.5]
            \begin{feynman}
                \vertex at (-2.7,1.5) (A);
                \vertex at (-2.7,-1.5) (B);
                \vertex at (-0.7,1) (C);
                \vertex at (-0.2,-1) (D);
                \vertex at (1,1.5) (E);
                \vertex at (1,-1.5) (F);
                \diagram*{
                    (A) -- [fermion, edge label={\(e^+\)}] (C);
                    (B) -- [fermion, edge label'={\(e^-\)}] (D);
                    (C) -- [boson, edge label={\(\gamma\)}] (E);
                    (D) -- [boson, edge label={\(\gamma\)}] (F);
                    (D) -- [fermion, edge label'={\( \)}] (C);
                };
            \end{feynman}
        \end{tikzpicture} 
    \end{minipage}
    \caption{Feynman diagrams for $e^+ e^-$ annihilation \cite{martin}}
    \label{dig:e+e-}
\end{figure}

% \newpage
\printbibliography[title={References}]

\end{document}
